\paragraph{Prediction from recorded offers.}
As a follow-up we fitted equal-size, row-permutation-equivariant networks to the actual 2025 regular offer or the outside option. Both networks receive program selection score, score relative to the previous cutoff, previous demand per seat, current capacity, and a missing-history indicator. One also receives reported preference rank. Options are permuted independently in every applicant row, lifting any dependence on column position. Each setting uses three H200 fits, a fixed 12 epochs and average predicted probabilities. There are no 2026 labels in the inference files. Both models are evaluated at the same four quotas in the same 2026 cohort as the primary analysis. The outcome-blind construction found the same 1,328,381 eligible pairs.

With five programs, the rank-blind model retained 144,401 of the recorded regular offers (92.42\%) and placed 158,800 applicants. Giving the same architecture the submitted rank retained 155,657 offers (99.62\%) and placed 158,533. This is a gain of 11,256 covered historical offers and a net loss of 267 counterfactual placements. Relative to full-list DA, the rank-informed model caused fewer members of the original placed set to lose a place (364 versus 1,045). It also placed far fewer previously unmatched people (200 versus 1,148). The additional 681 retained recipients and 948 fewer entrants sum to the 267-placement loss. This distinction is important when the service both predicts historical offers and evaluates policy performance on previous recipients.
